% Part: methods
% Chapter: proofs
% Section: inference-patterns

\documentclass[../../../include/open-logic-section]{subfiles}

\begin{document}

\olfileid{mth}{prf}{pat}

\olsection{અનુમાનના નમૂનાઓ}

સાબિતીઓ વ્યક્તિગત અનુમાનોથી બને છે. આપણે અનુમાન કરીએ ત્યારે
સામાન્ય રીતે ``તેથી'', ``આથી'' કે ``એટલે'' જેવા શબ્દોથી તે
દર્શાવીએ છીએ. અનુમાન ઘણી વાર સાબિતીમાં અગાઉથી ઉપલબ્ધ એક કે બે
હકીકતો પર આધાર રાખે છે---તે કોઈ ધારેલી બાબત હોઈ શકે અથવા
અગાઉના અનુમાનથી મેળવેલો નિષ્કર્ષ હોઈ શકે. સ્પષ્ટતા માટે આપણે
આવી બાબતોને ચિહ્નિત કરી શકીએ અને નવા અનુમાનમાં અન્ય કયાં વિધાનો
વાપરીએ છીએ તે જણાવી શકીએ. ઘણી વાર અનુમાનમાં અગાઉની બાબતોમાંથી
નવો નિષ્કર્ષ \emph{શા માટે} નીકળે છે તેની સમજ પણ હોય છે.
સાબિતીમાં વારંવાર વપરાતા અનુમાનના કેટલાક સામાન્ય નમૂનાઓ નીચે
જોઈશું. ગાણિતિક અનુમાનથી સાબિતી જેવા કેટલાક નમૂનાઓ વધુ વિસ્તૃત
છે (અને તેમની ચર્ચા પછી થશે).

અનુમાનનો એક નમૂનો આપણે જોઈ ચૂક્યા છીએ: વ્યાખ્યા ઉકેલવી અથવા
લાગુ પાડવી. વ્યાખ્યા ઉકેલીએ ત્યારે વ્યાખ્યેય પદ ધરાવતી બાબતને
વ્યાખ્યાકારક વિધાનથી ફરી કહીએ છીએ. દાખલા તરીકે, ધારો કે સાબિતીમાં
આપણે $D = E$ પહેલેથી સ્થાપિત કર્યું છે (a). હવે ગણોની સમાનતા
માટેની $=$ની વ્યાખ્યા લાગુ પાડીને કહી શકીએ: ``આથી, (a) અને
વ્યાખ્યા મુજબ,~$D$નું દરેક !!{element}~$E$નું !!a{element} છે
અને ઊલટું પણ છે.''

થોડું ગૂંચવણભર્યું એ છે કે અનુમાન કરતી ઘડીએ તેનો આધાર ઘણી વાર
લખતા નથી, પરંતુ પહેલાં લખીએ છીએ. ધારો કે $D = E$ હજી સાબિત
કર્યું નથી, પણ કરવું છે. જો $D = E$ આપણું લક્ષ્ય હોય, તો
વ્યાખ્યા લાગુ પાડીને લક્ષ્યને ફરી કહી શકીએ: $D = E$ સાબિત
કરવા~$D$નું દરેક !!{element}~$E$નું !!a{element} છે અને
ઊલટું પણ છે તે સાબિત કરવું પડે. તેથી સાબિતીનું સ્વરૂપ આવું હશે:
(a)~$D$નું દરેક !!{element}~$E$નું !!a{element} છે તે સાબિત
કરો; (b)~$E$નું દરેક !!{element}~$D$નું !!a{element} છે તે
સાબિત કરો; (c) તેથી (a), (b) અને $=$ની વ્યાખ્યા પરથી $D = E$.
પરંતુ સામાન્ય રીતે આપણે આવું લખતા નથી. તેના બદલે આવું લખીએ:
\begin{quote}
આપણે $D = E$ બતાવવું છે. $=$ની વ્યાખ્યા મુજબ તેનો અર્થ
એ છે કે~$D$નું દરેક !!{element}~$E$નું !!a{element} છે અને
ઊલટું પણ છે તે બતાવવું.

(a) \dots (~$D$નું દરેક !!{element}~$E$નું !!a{element} છે
તેની સાબિતી) \dots

(b) \dots (~$E$નું દરેક !!{element}~$D$નું !!a{element} છે
તેની સાબિતી) \dots
\end{quote}

\subsection{સંયોજનનો ઉપયોગ}

કદાચ અનુમાનનો સૌથી સરળ નમૂનો સંયોજનના બે ભાગમાંથી એકને
નિષ્કર્ષ તરીકે મેળવવાનો છે. એટલે, જો આપણે $p$ અને~$q$ ધાર્યાં
હોય કે અગાઉ સાબિત કર્યાં હોય, તો~$p$ (અને~$q$ પણ) અનુમાની
શકીએ. આ એટલું મૂળભૂત અનુમાન છે કે તેનો ઘણી વાર ઉલ્લેખ પણ
થતો નથી. દાખલા તરીકે, $D = E$ની વ્યાખ્યા ઉકેલ્યા પછી આપણે
સ્થાપિત કર્યું કે~$D$નું દરેક !!{element}~$E$નું !!a{element}
છે અને ઊલટું પણ છે. તેમાંથી~$E$નું દરેક !!{element}~$D$નું
!!a{element} છે એવો નિષ્કર્ષ કાઢી શકાય (આ ``ઊલટું પણ'' વાળો ભાગ છે).

\subsection{સંયોજન સાબિત કરવું}

ક્યારેક તમને સાબિત કરવા અપાયેલું વિધાન સંયોજનના સ્વરૂપનું હોય:
``$p$ અને $q$ સાબિત કરો.'' એ વખતે ફક્ત બે કામ કરવાનાં છે:
$p$ સાબિત કરો અને પછી $q$ સાબિત કરો. સ્પષ્ટતા માટે સાબિતીને
બે ભાગમાં વહેંચીને તેમના નામ આપી શકો. શરૂઆતની નોંધમાં
પાનાની ટોચે ``(1) $p$ સાબિત કરો'' અને પાનાની વચ્ચે
``(2) $q$ સાબિત કરો'' લખી શકો. (અલબત્ત, તમને સંયોજન
સાબિત કરવા સીધું ન કહેવાયું હોય, છતાં તમારી સાબિતીમાં તેની
જરૂર પડી શકે. દાખલા તરીકે, $D = E$ સાબિત કરવાનું હોય તો
$=$ની વ્યાખ્યા ઉકેલ્યા પછી તમારે સાબિત કરવાનું થાય છે કે
~$D$નું દરેક !!{element}~$E$નું !!a{element} છે
\emph{અને}~$E$નું દરેક !!{element}~$D$નું !!a{element} છે.)

\subsection{વિયોજન સાબિત કરવું}

સાબિત કરવાનું વિધાન વિયોજનના સ્વરૂપનું હોય (એટલે ``$p$ અથવા
$q$''ના સ્વરૂપનું વિધાન હોય), તો બે વિકલ્પમાંથી એક સાચો છે તે
બતાવવું પૂરતું છે. પરંતુ તમારા પ્રમેયની પૂર્વધારણામાંથી કોઈ એક
વિકલ્પ સીધો નીકળી આવે એવું લગભગ ક્યારેય બનતું નથી. વધુ વખત
પૂર્વધારણાઓ પોતે વિયોજક હોય છે, અથવા તમે બતાવતા હો છો કે અમુક
જાતની બધી વસ્તુઓમાં બેમાંથી એક ગુણધર્મ હોય છે---પણ કેટલીક
વસ્તુઓમાં પહેલો અને બીજીમાં બીજો ગુણધર્મ હોય છે. આવી સ્થિતિમાં
કિસ્સાવાર સાબિતી ઉપયોગી થાય છે (નીચે જુઓ).

\subsection{શરતી સાબિતી}

તમને મળતા ઘણાં પ્રમેયો શરતી સ્વરૂપના હશે (એટલે, $p$ સાચું હોય
તો $q$ પણ સાચું છે તે બતાવવાનું હોય). આવા કિસ્સાઓમાં સાબિતી
ગોઠવવી સરળ છે---શરતી વિધાનનો પૂર્વપક્ષ (અહીં $p$) ધારો અને
તેમાંથી પરિણામ~$q$ સાબિત કરો. એટલે તમારું પ્રમેય ``જો $p$ તો
$q$'' કહે, તો સાબિતીની શરૂઆત ``$p$ ધારો''થી કરો અને અંતે
~$q$ સાબિત થવું જોઈએ.

શરતી વિધાનો જુદી રીતે કહેવાય છે. ``જો $p$ તો $q$''ને બદલે
પ્રમેય ``$p$ ફક્ત ત્યારે જ જો $q$'', ``જો $p$ તો $q$'' અથવા
``$p$ હોય એ શરતે $q$'' એમ કહી શકે. આ બધાનો અર્થ એક જ છે:
$p$ ધારીને તેમાંથી~$q$ સાબિત કરવું. યાદ રાખો કે દ્વિશરતી
વિધાન (``$p$ ત્યારે અને ત્યારે જ જ્યારે $q$'', અંગ્રેજીમાં ``iff'') ખરેખર બે
શરતી વિધાનો છે: જો $p$ તો $q$, અને જો $q$ તો~$p$. તેથી
શરતી સાબિતીના બે દાખલા પૂરતા છે: પહેલા શરતી વિધાન માટે એક
અને બીજા માટે બીજો. પરંતુ ક્યારેક ``ત્યારે અને ત્યારે જ''નાં
વિધાનોની સાંકળથી પણ દ્વિશરતી વિધાન સાબિત કરી શકાય છે, જેમાં
``$p$''થી શરૂ કરીને ``$q$'' સુધી પહોંચીએ---પરંતુ ત્યારે
દરેક પગલું ખરેખર બંને દિશામાં ચાલે છે તેની ખાતરી કરવી પડે.

\subsection{સાર્વત્રિક દાવા}

સાર્વત્રિક દાવો વાપરવો સરળ છે: કોઈ બાબત બધી વસ્તુઓ માટે
સાચી હોય તો દરેક ચોક્કસ વસ્તુ માટે પણ સાચી છે. દાખલા તરીકે,
સાબિતીની પૂર્વધારણા $A \subseteq B$ હોય, તો $\subseteq$ની
વ્યાખ્યા ઉકેલતાં તેનો અર્થ એ છે કે દરેક $x \in A$ માટે
$x \in B$. તેથી $z \in A$ અગાઉથી જાણતા હો તો
$z \in B$ એવો નિષ્કર્ષ કાઢી શકો.

સાર્વત્રિક દાવો સાબિત કરવો થોડો અઘરો લાગશે. સામાન્ય રીતે આ
વિધાનોનું સ્વરૂપ ``જો $x$માં~$P$ ગુણધર્મ હોય તો તેમાં~$Q$
ગુણધર્મ પણ છે'' અથવા ``બધા $P$ એ $Q$ છે'' એવું હોય છે.
અલબત્ત, સ્વરૂપ હંમેશાં બરાબર એવું ન પણ હોય; તમને ચોક્કસ
શું સાબિત કરવાનું છે તે સમજવામાં થોડી કસરત જોઈએ. પરંતુ ઘણી
વાર કોઈ ગુણધર્મ ધરાવતી બધી વસ્તુઓમાં બીજો ગુણધર્મ પણ છે તે
સાબિત કરવું પડે છે.

સાર્વત્રિક દાવો સાબિત કરવા પ્રથમ ગુણધર્મ ધરાવતી વસ્તુઓને નામ
અથવા ચલ આપો અને પછી બતાવો કે તેમાં બીજો ગુણધર્મ પણ છે.
બીજી રીતે કહીએ તો, \emph{બધા}~$P$ માટે કંઈક સાબિત કરવા
\emph{મનસ્વી}~$P$ માટે તે સાબિત કરો. રજૂ કરાયેલું નામ કોઈ
પણ મનસ્વી~$P$ માટે છે. સામાન્ય રીતે મનસ્વી વસ્તુઓ માટે
એક અક્ષરનાં નામ વાપરીએ છીએ અને અક્ષરોની પરંપરા હોય છે:
જેમ કે પ્રાકૃતિક સંખ્યાઓ માટે $n$, !!{formula}s માટે $!A$,
ગણો માટે $A$, વિધેયો માટે $f$, વગેરે.

મુખ્ય કળા આખી સાબિતીમાં સામાન્યતા જાળવવાની છે. શરૂઆતમાં મનસ્વી
વસ્તુ (``$x$'')માં~$P$ ગુણધર્મ છે એમ ધારો અને ફક્ત વ્યાખ્યાઓ
અથવા માન્ય પૂર્વધારણાઓના આધારે બતાવો કે $x$માં~$Q$ ગુણધર્મ
પણ છે. તેમાં $P$ છે તે સિવાય $x$ ચોક્કસ શું છે તે નક્કી
કર્યું નથી; તેથી $P$ ધરાવતી દરેક વસ્તુમાં~$Q$ છે એમ કહી
શકો. ટૂંકમાં, $x$ એ~$P$ ગુણધર્મ ધરાવતી \emph{બધી} વસ્તુઓનું
પ્રતિનિધિત્વ કરે છે.

\begin{prop}
  બધા ગણો $A$ અને $B$ માટે $A \subseteq A \cup B$.
\end{prop}

\begin{proof}
  $A$ અને $B$ મનસ્વી ગણો લો. આપણે $A
  \subseteq A \cup B$ બતાવવું છે. $\subseteq$ની વ્યાખ્યા મુજબ
  તેનો અર્થ છે: દરેક $x$ માટે, જો $x \in A$ તો $x \in A \cup B$.
  તેથી~$x \in A$ને~$A$નું મનસ્વી !!{element} લો. આપણે $x \in A
  \cup B$ બતાવવું છે. $x \in A$ હોવાથી $x \in A$ અથવા $x \in B$.
  આથી $x \in \Setabs{x}{x \in A \lor x \in B}$. પરંતુ $\cup
  $ની વ્યાખ્યા મુજબ તેનો અર્થ $x \in A \cup B$ થાય છે.
\end{proof}

\subsection{કિસ્સાવાર સાબિતી}

ધારો કે પૂર્વધારણા અથવા અગાઉના નિષ્કર્ષ તરીકે તમારી પાસે વિયોજન
છે---તમે $p$ અથવા $q$ સાચું છે એમ ધાર્યું કે સાબિત કર્યું છે.
તમારે $r$ સાબિત કરવું છે. તે બે પગલાંમાં કરો: પહેલાં $p$ સાચું
છે એમ ધારીને~$r$ સાબિત કરો; પછી $q$ સાચું છે એમ ધારીને ફરી
~$r$ સાબિત કરો. આ રીત ચાલે છે કારણ કે બે વિકલ્પમાંથી એક
સાચો છે એવું આપણે ધારીએ છીએ અથવા જાણીએ છીએ. બંને પગલાં
બતાવે છે કે દરેક વિકલ્પ~$r$ના સત્ય માટે પૂરતો છે. (બંને
વિકલ્પ સાચા હોય તો $r$ સાચું છે તેનાં એક નહીં પણ બે કારણો
મળે છે. $p$ અને~$q$ બંને ધારીને $r$ અલગથી સાબિત કરવાની
જરૂર નથી.) આપણે શું કરીએ છીએ તે જણાવવા ``કિસ્સા અલગ
પાડીએ'' એમ કહીએ છીએ. દાખલા તરીકે, ધારો કે $x \in B \cup C$
છે. $B \cup C$ની વ્યાખ્યા $\Setabs{x}{x \in B \text{ અથવા } x \in C}$
છે. એટલે, વ્યાખ્યા મુજબ $x \in B$ અથવા $x \in C$. તેમાંથી
$x \in A$ સાબિત કરવું હોય, તો પહેલાં $x \in B$ ધારીને આ
પૂર્વધારણાથી $x \in A$ સાબિત કરીએ; પછી $x \in C$ ધારીને
ફરી તેમાંથી $x \in A$ સાબિત કરીએ. પૂર્વધારણા નીચે ``કિસ્સા
અલગ પાડીએ'' લખો; તેની નીચે ``કિસ્સો (1): $x \in B$'' અને
પાનાની વચ્ચે ``કિસ્સો (2): $x \in C$'' લખો. પછી પાનાના
ઉપરના અને નીચેના અર્ધભાગમાં સાબિતીઓ પૂરી કરો.

જે વિધાન સાબિત કરવાનું હોય તે પોતે વિયોજક હોય ત્યારે કિસ્સાવાર
સાબિતી ખાસ ઉપયોગી છે. આ સરળ ઉદાહરણ જુઓ:

\begin{prop}
ધારો કે $B \subseteq D$ અને $C \subseteq E$. તો $B \cup C \subseteq
D \cup E$.
\end{prop}

\begin{proof}
  (a) $B \subseteq D$ અને (b) $C \subseteq E$ ધારો. વ્યાખ્યા
  મુજબ, દરેક $x \in B$ માટે $x \in D$ પણ છે (c), અને દરેક
  $x \in C$ માટે $x \in E$ પણ છે (d).\sourcecorrection{OLMD-155}{મૂળમાં
  (c) અને (d)ના સભ્યતા-વિધાનમાં $x$ છૂટી ગયું છે; અહીં તે
  સ્પષ્ટ કર્યું છે.} $B \cup C \subseteq D \cup E$ બતાવવા
  આપણે બતાવવું છે કે જો $x \in B \cup C$ તો $x \in D \cup E$
  ($\subseteq$ની વ્યાખ્યા મુજબ). $x \in B \cup C$ ત્યારે
  અને ત્યારે જ જ્યારે $x \in B$ અથવા $x \in
  C$ ($\cup$ની વ્યાખ્યા મુજબ). એ જ રીતે, $x \in D \cup E$
  ત્યારે અને ત્યારે જ જ્યારે $x \in D$ અથવા $x \in E$.
  તેથી દરેક $x$ માટે બતાવવાનું છે કે જો $x \in B$ અથવા
  $x \in C$, તો $x \in D$ અથવા $x \in E$.

  \begin{quote}
  અત્યાર સુધી આપણે ફક્ત વ્યાખ્યાઓ ઉકેલી છે!{} આપણે
  પ્રતિજ્ઞાપનને $\subseteq$ અને $\cup$ વિના ફરી કહ્યું છે;
  હવે સાર્વત્રિક શરતી દાવો સાબિત કરવો છે. ઉપરની ચર્ચા મુજબ,
  આ માટે એવું $x$ ધારો જેના વિશે ``જો'' વાળો ભાગ સાચો છે
  અને પછી બતાવો કે ``તો'' વાળો ભાગ પણ સાચો છે. એટલે,
  $x
  \in B$ અથવા $x \in C$ ધારીને $x \in D$ અથવા $x \in
  E$ બતાવીશું.\footnote{આ ફકરો આપણે શું કરીએ છીએ તે જ સમજાવે
    છે---તે સાબિતીનો ભાગ નથી, અને તમે તમારી સાબિતી લખો ત્યારે
    આટલી વિગતમાં જવાની જરૂર નથી.}
  \end{quote}

  ધારો કે $x \in B$ અથવા $x \in C$. આપણે $x \in D$
  અથવા $x \in E$ બતાવવું છે. કિસ્સા અલગ પાડીએ.

  કિસ્સો 1: $x \in B$. (c) પરથી $x \in D$. તેથી $x \in D$ અથવા
  $x \in
  E$. (અહીં અગાઉના ઉપવિભાગમાં ચર્ચાયેલું અનુમાન કર્યું છે!)

  કિસ્સો 2: $x \in C$. (d) પરથી $x \in E$. તેથી $x \in D$ અથવા $x \in E$.
 \end{proof}

\subsection{અસ્તિત્વનો દાવો સાબિત કરવો}

અસ્તિત્વનો દાવો સાબિત કરવા કહેવાય ત્યારે પ્રશ્ન સામાન્ય રીતે
``એવો~$x$ છે કે $\dots x \dots$ તે સાબિત કરો'' એ સ્વરૂપનો
હશે; એટલે ``$\dots x
\dots$''થી વર્ણવાયેલો ગુણધર્મ ધરાવતી કોઈ વસ્તુ છે તે બતાવવાનું
હશે. એ માટે યોગ્ય વસ્તુ શોધીને બતાવવું પડે કે તેમાં જરૂરી
ગુણધર્મ છે. સાંભળવામાં સીધું લાગે છે, પણ આવી સાબિતી મુશ્કેલ
હોઈ શકે. સામાન્ય રીતે વસ્તુનું \emph{નિર્માણ} અથવા તેની
\emph{વ્યાખ્યા} આપવી પડે અને પછી સાબિત કરવું પડે કે એમ
વ્યાખ્યાયિત વસ્તુમાં જરૂરી ગુણધર્મ છે. યોગ્ય વસ્તુ શોધવી
મુશ્કેલ હોઈ શકે; તેમાં જરૂરી ગુણધર્મ છે તે સાબિત કરવું મુશ્કેલ
હોઈ શકે; અને ક્યારેક તમે ખરેખર કોઈ વસ્તુ વ્યાખ્યાયિત કરી છે
એ જ બતાવવું મુશ્કેલ હોય છે!{}

સામાન્ય રીતે તમે વસ્તુ નિર્દેશીને લખો, જેમ કે ``$x$ને
\dots લો'' (અહીં \dots{}થી કઈ વસ્તુ મનમાં છે તે કહો), જરૂર
હોય તો $\dots$ ખરેખર અસ્તિત્વ ધરાવતી વસ્તુ વર્ણવે છે તે
સાબિત કરો, અને પછી બતાવો કે $x$માં~$Q$ ગુણધર્મ છે. એક
સરળ ઉદાહરણ જુઓ.

\begin{prop}
  ધારો કે $x \in B$. તો એવો~$A$ છે કે $A \subseteq
  B$ અને $A \neq \emptyset$.
\end{prop}

\begin{proof}
  $x \in B$ ધારો. $A = \{x\}$ લો.
  \begin{quote}
    અહીં આપણે તેના !!{element}sની યાદી આપીને ગણ~$A$ની
    વ્યાખ્યા કરી છે. $x$ એક વસ્તુ છે એમ ધાર્યું છે અને તેના
    !!{element}sની યાદીથી ગણ હંમેશાં બનાવી શકાય છે; તેથી અહીં
    ગણ~$A$ની સફળતાપૂર્વક વ્યાખ્યા થઈ છે તે અલગથી બતાવવું
    પડતું નથી. પરંતુ પ્રતિજ્ઞાપન માગે છે તે ગુણધર્મો $A$માં
    છે એવું હજી બતાવવું પડે. તેના વિના સાબિતી પૂરી નથી!{}
  \end{quote}
  $x \in A$ હોવાથી $A \neq \emptyset$.
  \begin{quote}
    અહીં $A$ની $\{x\}$ તરીકેની વ્યાખ્યા અને $x \in \{x\}$
    તથા $x \notin \emptyset$ની સ્પષ્ટ હકીકતો વપરાય છે.
  \end{quote}
  $\{x\}$નું એકમાત્ર !!{element} $x$ છે અને $x \in B$ છે;
  તેથી~$A$નું દરેક !!{element}~$B$નું !!a{element} પણ છે.
  $\subseteq$ની વ્યાખ્યા મુજબ $A \subseteq B$.
\end{proof}

\subsection{અસ્તિત્વના દાવાનો ઉપયોગ}

ધારો કે અસ્તિત્વનો કોઈ દાવો સાચો છે તે તમે જાણો છો (કાં તો
તમે તેને સાબિત કર્યો છે અથવા તે વાપરી શકાય એવી પૂર્વધારણા છે),
જેમ કે ``કોઈ~$x$ માટે $x \in A$'' અથવા ``કોઈ $x \in A$
છે.'' સાબિતીમાં તેનો ઉપયોગ કરવો હોય તો પૂર્વધારણા મુજબ
અસ્તિત્વ ધરાવતી વસ્તુઓમાંથી એકને નામ આપ્યું છે એમ માની શકો.
$A$માં ઓછામાં ઓછી એક વસ્તુ છે, તેથી તે નામ જેનો ઉલ્લેખ
કરી શકે એવી વસ્તુઓ છે. તમે કદાચ તેમાંથી કોઈ ચોક્કસ વસ્તુ
પસંદ કરીને વધુ વર્ણવી ન શકો (તે $\in A$ છે એ સિવાય).
પરંતુ સાબિતી માટે તમે એક પસંદ કરી છે એમ માનીને તેને નામ
આપી શકો. અગાઉ વપરાયેલું અથવા પૂર્વધારણાઓમાં આવેલું નામ
ફરી ન વાપરવું અગત્યનું છે; નહીં તો ભૂલ થઈ શકે. સાબિતીમાં
``કોઈ $x$ માટે $x \in A$'' પરથી ``ધારો કે $a
\in A$'' તરફ જઈને આ દર્શાવો. હવે તમે~$a$ વિશે વિચાર કરી,
બીજી પૂર્વધારણાઓ વાપરીને વગેરે અંતે નિષ્કર્ષ $p$ સુધી પહોંચો.
જો $p$માં હવે~$a$નો ઉલ્લેખ ન હોય, તો $p$ ખાસ પસંદ કરેલા
$a \in A$ સાક્ષી પર આધાર રાખતું નથી; તે ``કોઈ $x$ માટે
$x \in A$'' એવી પૂર્વધારણાથી નીકળે છે.\sourcecorrection{OLMD-156}{મૂળમાં
નિષ્કર્ષને $a \in A$ની પૂર્વધારણાથી જ સ્વતંત્ર કહેવામાં આવ્યો છે.
અહીં ફક્ત ખાસ પસંદ કરેલા સાક્ષીથી સ્વતંત્રતા દર્શાવી છે;
અસ્તિત્વની પૂર્વધારણા હજી વપરાય છે.}

\begin{prop}
જો $A \neq \emptyset$, તો $A \cup B \neq \emptyset$.
\end{prop}

\begin{proof}
  $A \neq \emptyset$ ધારો. તેથી કોઈ $x$ માટે $x \in A$.
  \begin{quote}
    અહીં પહેલાં આપણે પ્રતિજ્ઞાપનની પૂર્વધારણા ફરી કહી છે. આ
    પૂર્વધારણા, એટલે $A \neq \emptyset$, અસ્તિત્વનો દાવો
    છુપાવે છે; કેટલીક વ્યાખ્યાઓ ઉકેલીએ ત્યારે તે મળે છે. $=$ની
    વ્યાખ્યા કહે છે કે $A = \emptyset$ ત્યારે અને ત્યારે જ જ્યારે
    દરેક $x \in A$ માટે $\in \emptyset$ પણ છે અને દરેક $x \in
    \emptyset$ માટે $\in A$ પણ છે. બંને બાજુનો નિષેધ કરતાં,
    $A \neq
    \emptyset$ ત્યારે અને ત્યારે જ જ્યારે કોઈ $x \in A$ માટે
    $\notin \emptyset$ છે અથવા કોઈ $x \in \emptyset$ માટે
    $\notin A$ છે. કોઈ વસ્તુ~$\in \emptyset$ નથી; તેથી બીજો
    વિકલ્પ ક્યારેય સાચો નથી, અને ``$x \in A$ અને $x
    \notin \emptyset$'' ફક્ત $x \in A$ બને છે. તેથી $A \neq
    \emptyset$ ત્યારે અને ત્યારે જ જ્યારે કોઈ $x$ માટે
    $x \in A$.\sourcecorrection{OLMD-157}{મૂળના આ વાક્યમાં
    ખોટી રીતે x અસમાન શૂન્યગણ લખાયું છે; સંદર્ભ અને અગાઉના
    પગલાં મુજબ A અસમાન શૂન્યગણ હોવું જોઈએ.} આ અસ્તિત્વનો
    દાવો છે. હવે~$A$ના !!{element}sમાંથી એકને નામ આપીને આ
    દાવો વાપરીએ:
  \end{quote}
  $a \in A$ લો.
  \begin{quote}
    હવે $\in A$ એવી એક વસ્તુને નામ આપ્યું છે. આગળ~$a$
    વિશે વિચારીશું, પરંતુ ફક્ત $a \in A$ જ ધારીશું, બીજું
    કંઈ નહીં:
  \end{quote}
  $a \in A$ હોવાથી $\cup$ની વ્યાખ્યા મુજબ $a \in A \cup B$.
  તેથી કોઈ $x$ માટે $x \in A \cup B$, એટલે $A \cup B \neq \emptyset$.
  \begin{quote}
    છેલ્લા પગલામાં ``$a \in A \cup B$'' પરથી ``કોઈ
    $x$ માટે $x \in A \cup B$'' મેળવ્યું. તેમાં હવે $a$નો
    ઉલ્લેખ નથી; એટલે ``કોઈ $x$ માટે $x \in A \cup B$''
    ફક્ત ``કોઈ $x$ માટે $x \in A$'' પરથી નીકળે છે. તેનો
    અર્થ $A \cup B \neq
    \emptyset$ છે.
  \end{quote}
\end{proof}

આપણે કર્યું તેમ ``$x$'' જેવા બંધિત ચલોને $a$ જેવાં કામચલાઉ
નામોથી જુદાં રાખવાનો અભ્યાસ સારો હોઈ શકે. પરંતુ વ્યવહારમાં
ઘણી વાર આમ કરતા નથી અને $x$ જ વાપરીએ છીએ, જેમ કે:
\begin{quote}
ધારો કે $A \neq \emptyset$, એટલે કોઈ $x \in A$ છે. $\cup$ની
વ્યાખ્યા મુજબ $x \in A \cup B$. તેથી $A \cup B \neq \emptyset$.
\end{quote}
પરંતુ આમ કરતાં જુદા અસ્તિત્વના દાવાઓ માટે જુદા $x$ અને $y$
વાપરો છો તેની ખાસ કાળજી રાખવી પડે. દાખલા તરીકે, નીચેનું
``જો $A \neq
\emptyset$ અને $B \neq \emptyset$ તો $A \cap B \neq \emptyset$''
સાબિત કરવાનું \emph{સાચું નથી} (કારણ કે એ વિધાન જ સાચું નથી):
\begin{quote}
ધારો કે $A \neq \emptyset$ અને $B \neq \emptyset$. તેથી કોઈ $x$ માટે $x
\in A$ અને કોઈ $x$ માટે $x \in B$. $x \in A$ અને $x \in
B$ હોવાથી $\cap$ની વ્યાખ્યા મુજબ $x \in A \cap B$. તેથી $A \cap B \neq
\emptyset$.
\end{quote}
ખોટું પગલું ક્યાં આવે છે અને પરિણામ કેમ સાચું નથી તે તમે
શોધી શકો છો?

\end{document}
